% Fragmento público generado desde propietarios íntegros.
% Residencia: 52.13.3.
% Fuente: ../incoming/radion_monografia_20260827/manuscrito/sections/13_sintesis.tex:185-247
\subsubsection{Negative controls of the rational system}

The construction is refutable at specific points:
\begin{enumerate}
\item If \(S_E\notin(1,2)\), the APP quotient is not one and the closure changes
cell.
\item If the incidencia activates no produces \(N_6=90\), o if no there are two sheets,
\(20/81\) ceases of seguirse.
\item If the verifier receives \(180/\pi\) or \(\epsR\) before constructing the
operands, the generation is circular.
\item If subtraction by triads does not stabilize compatible prefixes, the
proposed coinductive section does not exist.
\item If \(|C^*|\ge A\), the positive ellipse ceases to exist.
\item If \(a_Sb_S\ne2\hret\), the area theorem fails.
\item If \(C^*=0\), \(d=0\), \(\eta=0\), \(\tau=i\) and
\(j=1728\) must hold.
\item If the Klein area of the subdisk does not follow \eqref{eq:klein-area-R}, the
infinite-depth thesis fails.
\item If \(F_{\mathrm{spec}}=\epsR^\circ\) is asserted without an independent counterterm, the
numerical difference \eqref{eq:chiR} refutes it.
\item Producing \(12:-1\) requires preserving, alongside the cardinalities
\(90/120\), the twelve dodecaphasic sectors, the unique oriented seam, and
the reader's action on these generators. Changing any multiplicity
must modify the image \(12S_{90}-S_{120}\). For the unaltered fundamental
chain, the coefficient of \((1-q)^{-1}\) must yield \(1/24\).
\item If the function \(j\) is used to select the charge sign, the
invariance \eqref{eq:j-invariance} is lost and a type error is incurred.
\end{enumerate}

\subsubsection{Joint physical realization of the radion}

The final image is not a corrected circle. It is an architecture of
publications:
\[
\begin{aligned}
&\text{APP:}&&\text{unit, remainder, carry and fiber};\\
&\text{TRIT:}&&\text{regime, orientation and sheets};\\
&\text{TPK:}&&\text{phase, return, dodecaphase and memory};\\
&\text{ellipse:}&&\text{action and anisotropy};\\
&\text{Klein:}&&\text{interior depth};\\
&\text{helix:}&&\text{biography of the return};\\
&\text{MD:}&&\text{clock, energy, charge, field and causality}.
\end{aligned}
\]

The radion equation brings together the visible face:
\BigRadionEquation
The ellipse brings together the action face:
\[
\mathcal A(1)=\hret,\qquad \mathcal A(2\pi)=\hfull.
\]
The Klein metric unites the depth:
\[
K=-1,\qquad \operatorname{Area}_K=\infty.
\]
Holonomy gathers memory:
\[
g(k+9)=g(k),\qquad B_{k+9}\ne B_k.
\]

\begin{thesisbox}
The circle is the angular shadow; the ellipse is the body of action; Klein measures
the interior; the helix preserves the biography. The radion is the unit that
allows the four statements to be read on a single state without reducing one
to another.
\end{thesisbox}
% Fuente: ../incoming/radion_monografia_20260827/manuscrito/sections/B_valores_y_verificacion.tex:1-105
\subsubsection{Canonical values, stability, and reproducible verification}

\paragraph{Values of the generated internal invariants}

\begin{longtable}{p{.27\textwidth} p{.61\textwidth}}
\toprule
Object & Value used in the coinductive profile\\
\midrule
\endfirsthead
\toprule
Object & Value used in the coinductive profile\\
\midrule
\endhead
\(\pih\) & \(\begin{aligned}3.1415926535897932384626433832795028841\\[-2pt]971693993751058209749445923\ldots\end{aligned}\)\\
\(e_{\HMT}\) & coinductive evaluation of the Euler reader, without an external table\\
\(\alphah\) & \(\begin{aligned}0.0072973525692838009972851054723806629\\[-2pt]995641725396427091854046825\ldots\end{aligned}\)\\
\(A\) & \(\begin{aligned}7.2973525692838009972851054723806629995\\[-2pt]641725396427\ldots{}^\circ\end{aligned}\)\\
\(C^*\) & \(2.1237383389686457242619513803933607937\ldots{}^\circ\)\\
\(\eta\) & \(0.299689683921630799684926562863927\ldots\)\\
\(\hret\) & \(1.05457181691503906113369058472430\ldots\times10^{-34}U_S\)\\
\bottomrule
\end{longtable}

\paragraph{Quantitative convergence with depth}

\begin{center}
\small
\begin{tabular}{r l l}
\toprule
Depth & \(\epsR^\circ\) & reading\\
\midrule
12 & \(5.1681282186551375597\times10^{-8}\) & preasymptotic\\
18 & \(5.1383017764316644535\times10^{-8}\) & 7 cifras estables\\
24 & \(5.1383016758575394560\times10^{-8}\) & 14 cifras estables\\
30 & \(5.1383016758575282072\times10^{-8}\) & 20 cifras estables\\
36 & \(5.138301675857528207168517\times10^{-8}\) & perfil largo\\
45 & \(5.13830167585752820716851646226459\times10^{-8}\) & width \(<4.81\times10^{-130}\)\\
\bottomrule
\end{tabular}
\end{center}

The depth returns \(9,18,27,36,45\) preserve the phase \(8\) and
increase the memory \(0,1,2,3,4\). The ratio between consecutive widths tends
to
\[
3^{-54}=1.7196982334549856127610399316004871\ldots\times10^{-26}.
\]

\paragraph{Structural Sensitivity Analysis by Controlled Suppression}

\begin{longtable}{p{.46\textwidth} p{.34\textwidth}}
\toprule
Operation & Output or unit difference\\
\midrule
\endfirsthead
\toprule
Operation & Output or unit difference\\
\midrule
\endhead
Canonical operator & \(8.96802822044562981999\times10^{-10}\)\\
A sheet, \(\rho=90/729\) & \(-1.3727056615960678\times10^{-5}\)\\
Phase \(\theta_1=20^\circ\) & \(+5.2359877649510169\times10^{-1}\)\\
Phase \(\theta_3=80^\circ\) & \(-5.2359877470149605\times10^{-1}\)\\
Without channel \(e\) & \(1.8065350748904653\times10^{-3}\)\\
Historical Torsion \(\Delta_4(A,C^*)\) & \(1.0517084608768816\times10^{-9}\)\\
Harmonic operator \(M_4\) minus canonical & \(1.0525500222895070\times10^{-17}\)\\
Hexagonal Cola less canonical & \(2.0283936357433839\times10^{-12}\)\\
\bottomrule
\end{longtable}

The separation of scales reveals structural sensitivity: closure depends
on the phase, the two sheets, the channels \(\pi/e\), global torsion, and carry.
No ablation reproduces the value.

\paragraph{Reproducible Computational Certificates}

The work uses the following executable proprietors:
\begin{enumerate}
\item \sourcepath{output/RADION_TWOWAY_NONADICO_20260827/verificar_emergencia_radion_por_profundidad.py};
\item \sourcepath{output/RADION_TWOWAY_NONADICO_20260827/verificar_radion_helice_critica_tpk.py};
\item \sourcepath{output/RADION_TWOWAY_NONADICO_20260827/verificar_operador_tpk_target_free_profundidad.py};
\item the consolidated verifier for this monograph, located at
\sourcepath{certificados/verificar_monografia_radion.py}.
\end{enumerate}
The expected outputs include
\[
\texttt{PASS\_EMERGENCIA\_RADION\_POR\_PROFUNDIDAD},
\]
\[
\texttt{PASS\_RADION\_HELICE\_CRITICA\_TPK},
\qquad
\texttt{PASS\_MONOGRAFIA\_RADION\_HMT\_MD}.
\]

\paragraph{Distinction between computational precision and mathematical accuracy}

The fact that a program displays zero to 120 or 200 digits does not by itself turn a residual identity into a prediction. Mathematical exactness derives from \eqref{eq:cierre-unitario}; arbitrary precision makes it possible to evaluate its operands and verify convergence. Non-circularity derives from the order of construction and the prohibition on target values, not from the number of digits.

Decimal computation may leave residuals of order \(10^{-218}\) when combining
branches evaluated at different precisions. That number is implementation
rounding, not an additional physical correction.
% Fuente: ../incoming/radion_monografia_20260827/manuscrito/sections/D_historia_fuentes.tex:1-169
\subsubsection{Conceptual Genealogy and Primary Sources}

\paragraph{Planck action}

Planck introduces \(h\) into the radiation law as the constant linking energy
and frequency. The decisive historical datum for this work is dimensional:
\(h\) is action. The radion does not alter that role; it explains why the action of
one complete turn and the action per unit of phase are related by \(2\pi\).

\paragraph{Bohr-Sommerfeld Angular Quantization}

Bohr uses \(h/2\pi\) in the quantization of angular momentum. Sommerfeld extends
quantization to action integrals and Keplerian orbits. Therefore, it is not
historically accurate to say that Dirac invented the reduced constant. His
contribution belongs to a genealogy in which angular action was already
present.

The HMT--MD novelty is different: it does not take \(h/2\pi\) as an external quotient,
but rather constructs an ellipse whose turn has area \(h\) and whose sector of one
radian has area \(\hbar\).

\paragraph{Born and Jordan Matrix Formalism and the Quantum Oscillator}

Born and Jordan place the variational principle, phases, and the harmonic
oscillator at the heart of matrix mechanics; the electromagnetic field is
decomposed into a collection of oscillators. The LC realization of the radion
continues that intuition with one additional datum: the form quotient
\(e^{-\eta}\) determines the relative impedance while product, frequency,
and action remain fixed.

\paragraph{Angular frequency and action in the Dirac formulation}

In 1925, Dirac writes frequencies in radians per unit time and uses
relations of the form
\[
xy-yx=\frac{ih}{2\pi}[x,y],
\qquad
\Delta E=\frac{h}{2\pi}\omega.
\]
The equation \(E=\hret\omega\) explicitly joins reduced action and
angular frequency. In the radion, that relation acquires a geometry: \(\hret\)
is the area swept by one unit of phase.

\paragraph{Heisenberg, Robertson, and Schrödinger Uncertainty Relations}

Heisenberg formulates the joint limitation of canonically
conjugate variables. Robertson generalizes the inequality; Schrödinger incorporates
covariance. The matrix \eqref{eq:covariance} belongs more directly to that
Robertson--Schrödinger prolongation than to Heisenberg's article alone.

None of these authors formulates the complete genealogical ellipse of this work.
The historically defensible assertion is that they provide the language in
which the action ellipse can be recognized as a minimal contour; the
construction of the pair \((A,C^*)\), the foci, Klein, the helix, and the carry is the
HMT--MD contribution.

\paragraph{Feynman Phase Functional}

The spacetime formulation assigns a phase to each path,
\(e^{iS/\hbar}\). The radion unit makes the correspondence literal:
incrementing the action by \(\hbar\) increments the phase by one radian. The sum over
paths recognizes, in integral language, the same principle that the ellipse
expresses as area.

\paragraph{Advanced and retarded propagations in Wheeler--Feynman}

Absorber theory constructs a temporally symmetric classical interaction
through combinations of retarded and advanced fields. The HMT
identity \(CU(t)C^{-1}=U(-t)\) and route inversion provide an
operator-theoretic antecedent. Physical promotion requires the map between routes
and propagators to preserve action, orientation, boundary, and symplectic form.

Absorber theory is not the same assertion as the Feynman--Stueckelberg
interpretation of antiparticles as temporally reversed propagation.
Both may recognize bidirectionality; their physical objects and
boundary conditions differ.

\paragraph{Hodge Duality and Separation with Respect to Hochschild Cohomology}

The operation relevant to this manuscript is the Hodge star \(\star^2=-I\) on two-forms in Lorentzian signature. No proprietor of Hochschild homology acting on the ellipse, the carry, or the sector \(90/120\) has been located in the active radion corpus. Introducing Hochschild by nominal similarity would degrade the typing. If a future prolongation constructs a deformation cocycle with its own domain and codomain, it must enter as a new result, not as a retroattribution.

\paragraph{Selected Primary Bibliography}

\begin{enumerate}
\renewcommand{\labelenumi}{[\arabic{enumi}]}
\item\hypertarget{radionbib:Planck1901}{}
M.~Planck,
\emph{On the Law of Energy Distribution in the Normal Spectrum},
Annalen der Physik 309 (1901), 553--563.
\url{https://doi.org/10.1002/andp.19013090310}

\item\hypertarget{radionbib:Bohr1913}{}
N.~Bohr,
\emph{On the Constitution of Atoms and Molecules},
Philosophical Magazine 26 (1913), pp. 1--25.
\url{https://doi.org/10.1080/14786441308634955}

\item\hypertarget{radionbib:Sommerfeld1916}{}
To.~Sommerfeld,
\emph{Zur Quantentheorie der Spektrallinien, I},
Annalen der Physik 356 (1916), 1--94.
\url{https://doi.org/10.1002/andp.19163561702}

\item\hypertarget{radionbib:BornJordan1925}{}
M.~Born and P.~Jordan,
\emph{Zur Quantenmechanik},
Zeitschrift fuer Physik 34 (1925), 858--888.
\url{https://doi.org/10.1007/BF01328531}

\item\hypertarget{radionbib:Dirac1925}{}
P.~A.~M.~Dirac,
\emph{The Fundamental Equations of Quantum Mechanics},
Proceedings of the Royal Society A 109 (1925), pp. 642--653.
\url{https://doi.org/10.1098/rspa.1925.0150}

\item\hypertarget{radionbib:Heisenberg1927}{}
W.~Heisenberg,
\emph{On the Perceptual Content of Quantum-Theoretical Kinematics and Mechanics},
Journal of Physics 43 (1927), 172--198.
\url{https://doi.org/10.1007/BF01397280}

\item\hypertarget{radionbib:Dirac1928}{}
P.~A.~M.~Dirac,
\emph{The Quantum Theory of the Electron},
Proceedings of the Royal Society, Series A 117 (1928), 610--624.
\url{https://doi.org/10.1098/rspa.1928.0023}

\item\hypertarget{radionbib:Robertson1929}{}
H.~P.~Robertson,
\emph{The Uncertainty Principle},
Physical Review 34 (1929), pp. 163--164.
\url{https://doi.org/10.1103/PhysRev.34.163}

\item\hypertarget{radionbib:Schrodinger1930}{}
E.~Schrodinger,
\emph{Zum Heisenbergschen Unscharfeprinzip},
Sitzungsberichte der Preussischen Akademie der Wissenschaften (1930), 296--303.

\item\hypertarget{radionbib:WheelerFeynman1945}{}
J.~A.~Wheeler and R.~P.~Feynman,
\emph{Interaction with the Absorber as the Mechanism of Radiation},
Reviews of Modern Physics 17 (1945), 157--181.
\url{https://doi.org/10.1103/RevModPhys.17.157}

\item\hypertarget{radionbib:Feynman1948}{}
R.~P.~Feynman,
\emph{Space-Time Approach to Non-Relativistic Quantum Mechanics},
Reviews of Modern Physics 20 (1948), pp. 367--387.
\url{https://doi.org/10.1103/RevModPhys.20.367}

\item\hypertarget{radionbib:Feynman1949}{}
R.~P.~Feynman,
\emph{The Theory of Positrons},
Physical Review 76 (1949), pp. 749--759.
\url{https://doi.org/10.1103/PhysRev.76.749}
\end{enumerate}

\begin{narrativebox}
The founders contain the components: action, angular frequency, conjugate
variables, oscillator, uncertainty, path phase, and time reversal.
The contribution of the radion is to bring them together in a generated object: an ellipse of
area \(h\), a sector of area \(\hbar\), an interior of infinite depth,
a sheet orientation, and an exact carry of compatibility.
\end{narrativebox}
% Fuente: ../incoming/radion_monografia_20260827/manuscrito/sections/E_conclusion.tex:1-107
\subsubsection{Final Synthesis: Angular unity with phase genealogy, action and memory}

The radian does not disappear. It becomes fully visible.

In its ordinary use, a radian is a ratio that need not remember how it was
produced. The arc and the radius cancel dimensionally; the unit appears
bare. That bareness is a virtue for trigonometry, but it is a forgetting
for a theory that seeks to explain the structure of the continuum and the origin
of action. The radion recovers what was forgotten without altering the published value.

The genealogy begins in APP, where an output preserves its value and also its sheet, quotient, remainder, carry, and boundary. TRIT introduces regime and orientation. TPK returns the phase without resetting the state. The nonadic connection constructs the coinductive section of \(\pi\), and the dodecaphase publishes the hinge of \(50^\circ\). The angular pair supplies a common mode \(A\) and an oriented opening \(C^*\). The incidence of six active phases produces \(90\), its two sheets over the ambient space \(729\) produce \(20/81\), and global torsion supplies the second section.

Then \(\varepsilon_R\) appears. Not as a decimal begged of the
closure, but as the difference that the structure was required to
transport:
\[
\epsone=\frac{20}{81}\Delta_4-(S_E-1).
\]
The carry publishes its digits; the depth certifies its limit; the angular
chart recognizes it as
\[
5.13830167585752820716851646226459474899\ldots
\times10^{-8}\,{}^\circ.
\]
The radian equation is closed because both sides are publications of the
same state:
\BigRadionEquation

The account does not end with the equality. The pair \((A,C^*)\) constructs a
positive ellipse. The product of its semiaxes fixes \(2\hret\); their quotient
fixes its shape. The area theorem proves that one phase unit sweeps
\(\hret\) and that one revolution encloses \(\hfull\). Division by
\(2\pi\) here acquires physical content: \(h\) is action per
revolution; \(\hbar\) is action per radion.

The foci cease to be ornaments in a figure. Their separation reconstructs
\(C^*\) in the spectral chart, and their Klein distance measures a finite
hyperbolic depth within a domain whose total hyperbolic area is
infinite. The same boundary contains a finite action. The apparent paradox
is the central thesis of the geometry: one measure closes the action; another opens
the depth.

The circle, therefore, is not false. It is the exact projection that preserves phase
and forgets anisotropy, interior, and memory. The ellipse preserves action and form.
The Klein surface preserves depth. The helix preserves return and biography. No
image replaces the others; all arise from the same state.

The Hilbert realization recognizes the ellipse as a contour of minimal covariance. The LC oscillator recognizes it as an exchange of electrical and magnetic reserves. The reader \(90/120\) publishes permeability, permittivity, impedance, and velocity. The six--eight incidence produces a causal screen; Hodge converts orientation into field duality; Lorentz converts field and charge into dynamics; Cartan--Holst receives the holonomy in gravitational geometry. The order matters: the HMT object selects; conventional physics realizes.

The word \emph{radion} is thereby justified. It does not mean that one focus is
a positive charge and the other a negative charge. It means that the phase
unit has a signed orientation prior to electrical publication:
\[
\mathcal A(\mathfrak r_+)=+\hret,
\qquad
\mathcal A(\mathfrak r_-)=-\hret.
\]
Charge appears when Dimensional Mechanics applies the transducer. The sign
is not added at the end; it is preserved from the sheet.

This also explains why the \(90/120\) sector appeared to conceal many
answers. Incidence, the dodecaphase extractor, the spectral tail, the Barbero
kernel, towers, and carry share a common origin. Precisely because they share an origin,
they can be compared. Precisely because they employ distinct operations, they must not
be merged. Typing does not diminish the unity of the system: it makes it intelligible.

% REMATE_LOCAL_RADION_CIERRE
\Needspace{28\baselineskip}
The most compact conclusion of this work can be written in four lines:
\[
\boxed{
\begin{aligned}
\text{phase:}\quad&\theta=1,\\
\text{action:}\quad&\mathcal A=\sigma\hret,\\
\text{memory:}\quad&H_9(n+9)=H_9(n)+(0,1),\\
\text{compatibility:}\quad&1=S_E-\Phi_T+\epsone.
\end{aligned}}
\]

The radion is that fourfold unit. The radian is its exact shadow in the
angular chart. The difference lies not in the value of \(\pi\) but in
how much structure we are willing to preserve when we say that a revolution
has closed.

\vspace{2\baselineskip}
\begin{center}
\begin{minipage}{.82\textwidth}
\centering\large\itshape\color{RadNavy}
The circle publishes phase.\\
The ellipse contains action.\\
Klein opens depth.\\
The helix remembers the path.\\[.7em]
The radion is the unit that does not force a choice among them.
\end{minipage}
\end{center}
